\section{Análisis exploratorio profundo}

El EDA se realiza sobre una fila por canción. Este cambio hace que cada estadística responda a la
unidad de negocio elegida y no a la cantidad de etiquetas de género asociadas.

\subsection{Distribución procesada de popularidad}

La media es 33,198, la mediana 33, la desviación estándar 20,578, los percentiles 25 y 75 son 19 y
49, y el percentil 95 es 67. El rango completo 0--100 es compatible con la escala documentada.

\begin{figure}[H]
\centering
\includegraphics[width=0.84\textwidth]{05_processed_popularity_distribution.png}
\caption{Popularidad de 89.740 canciones únicas. La concentración baja-media y la masa en cero son
relevantes para líneas base y evaluación por segmentos.}
\label{fig:popularidad-procesada}
\end{figure}

La Figura~\ref{fig:popularidad-procesada} importa porque define la dificultad del objetivo y
advierte que una métrica global puede ocultar errores distintos en canciones de baja y alta
popularidad. Los valores 0 y 100 no son errores solo por estar en los extremos. La limitación es que
la distribución describe una captura histórica sin fecha de referencia.

\subsection{Pearson y Spearman}

Pearson resume covariación lineal; Spearman calcula asociación monotónica sobre rangos y es menos
dependiente de una relación lineal exacta. La comparación evita concluir a partir de una sola forma.

\begin{table}[H]
\centering
\caption{Correlaciones con \code{popularity}, ordenadas por magnitud de Spearman.}
\label{tab:correlaciones-target}
\begin{tabular}{lrr}
\toprule
 & Pearson & Spearman \\
variable &  &  \\
\midrule
instrumentalness & -0.127 & -0.124 \\
loudness & 0.072 & 0.067 \\
speechiness & -0.047 & -0.067 \\
danceability & 0.064 & 0.055 \\
energy & 0.014 & -0.016 \\
duration\_ms & -0.023 & 0.013 \\
liveness & -0.014 & -0.012 \\
valence & -0.012 & -0.011 \\
acousticness & -0.039 & 0.010 \\
tempo & 0.007 & 0.008 \\
\bottomrule
\end{tabular}

\end{table}

Todas las magnitudes son pequeñas. \code{instrumentalness} es la señal negativa más visible;
\code{loudness} y \code{danceability} son positivas débiles. Diferencias de signo entre métodos,
como en \code{acousticness} y \code{duration_ms}, sugieren formas no lineales, empates o patrones
segmentados. La implicancia es evaluar modelos multivariables sin esperar gran poder explicativo de
un atributo aislado. La limitación central es que correlación no controla confusores ni implica
causalidad.

\subsection{Matriz de correlación y redundancia}

\begin{figure}[H]
\centering
\includegraphics[width=0.91\textwidth]{06_numeric_correlation_heatmap.png}
\caption{Matriz de Pearson para variables numéricas. Las celdas más intensas aparecen entre
predictores de audio, no con el objetivo.}
\label{fig:matriz-correlacion}
\end{figure}

La Figura~\ref{fig:matriz-correlacion} permite detectar redundancia: \code{energy}--\code{loudness}
alcanza 0,759; \code{energy}--\code{acousticness}, $-0{,}733$; y
\code{loudness}--\code{acousticness}, $-0{,}583$. Para regresión lineal, coeficientes asociados
pueden volverse inestables y difíciles de interpretar. Esto no obliga a retirar variables: permite
planificar regularización y diagnóstico. La matriz solo refleja relaciones pares y lineales.

\subsection{Relaciones bivariadas por cuantiles}

Un scatterplot con casi 90 mil puntos produciría sobreposición. Se dividen valores en cuantiles y se
representa la mediana de popularidad de cada intervalo, manteniendo todos los datos en el cálculo.

\begin{figure}[H]
\centering
\includegraphics[width=\textwidth]{07_binned_feature_popularity.png}
\caption{Popularidad mediana por cuantiles de cinco atributos. Las curvas permiten observar forma y
no solo un coeficiente global.}
\label{fig:cuantiles-popularidad}
\end{figure}

La Figura~\ref{fig:cuantiles-popularidad} muestra caída más consistente con
\code{instrumentalness}; \code{loudness} y \code{danceability} crecen en parte del rango, pero no de
forma ilimitada. \code{energy} y \code{acousticness} exhiben curvatura. Esto respalda términos no
lineales o modelos flexibles como hipótesis de modelamiento, no reglas creativas. Los cuantiles
resumen heterogeneidad y pueden ocultar subgrupos por género o época.

\subsection{Contenido explícito y tamaño de efecto}

\begin{figure}[H]
\centering
\includegraphics[width=0.72\textwidth]{08_explicit_popularity_boxplot.png}
\caption{Distribución de popularidad por clasificación explícita. El solapamiento entre grupos es
amplio pese a diferencias en los centros.}
\label{fig:explicit-popularidad}
\end{figure}

El grupo no explícito contiene 82.036 canciones y el explícito 7.704. Sus medias son 32,852 y
36,884; las medianas, 33 y 37. Para medir relevancia práctica se calculó $d$ de Cohen con desviación
estándar combinada \citep{cohen1988}.

\begin{table}[H]
\centering
\caption{Diferencias y tamaño de efecto de \code{explicit}.}
\label{tab:explicit-efecto}
\begin{tabular}{lr}
\toprule
 & Valor \\
\midrule
Canciones no explicitas & 82036.000 \\
Canciones explicitas & 7704.000 \\
Diferencia de medias & 4.032 \\
Diferencia de medianas & 4.000 \\
d de Cohen & 0.196 \\
\bottomrule
\end{tabular}

\end{table}

La diferencia de 4,03 puntos puede resultar estadísticamente detectable con una muestra grande,
pero $d=0{,}196$ indica separación estandarizada pequeña. La Figura~\ref{fig:explicit-popularidad}
confirma fuerte solapamiento. La asociación puede depender de género, público, recencia y
promoción; no prueba que hacer una canción explícita aumente popularidad.

\subsection{Géneros multietiqueta y sensibilidad}

Cada valor de \code{track_genres} se divide por \code{|}. Una canción puede participar en varias
etiquetas: las filas expandidas no constituyen canciones adicionales. El resumen cuenta IDs únicos
por etiqueta y exige soporte para evitar líderes de muestras mínimas.

\begin{figure}[H]
\centering
\includegraphics[width=0.82\textwidth]{09_genre_popularity_summary.png}
\caption{Géneros con mayor mediana de popularidad entre etiquetas con al menos 100 canciones.}
\label{fig:generos-popularidad}
\end{figure}

\begin{table}[H]
\centering
\caption{Sensibilidad del líder por mediana a distintos soportes mínimos.}
\label{tab:sensibilidad-generos}
\begin{tabular}{lrlr}
\toprule
 & Generos incluidos & Genero lider & Mediana lider \\
Soporte minimo &  &  &  \\
\midrule
100 & 114 & pop & 66.000 \\
200 & 114 & pop & 66.000 \\
500 & 114 & pop & 66.000 \\
\bottomrule
\end{tabular}

\end{table}

Los 114 géneros superan 100, 200 y 500 canciones, y \code{pop} mantiene mediana 66. Esto demuestra
estabilidad frente a esos umbrales en esta muestra, no superioridad universal. La Figura
\ref{fig:generos-popularidad} sirve para segmentar el dataset; su limitación es el balance artificial
del origen y la superposición de canciones entre etiquetas.

\subsection{Valores extremos}

El criterio IQR marca observaciones fuera de $[Q_1-1{,}5IQR, Q_3+1{,}5IQR]$. Es una regla de
identificación, no una prueba de error.

\begin{table}[H]
\centering
\caption{Rangos y recuentos fuera de IQR.}
\label{tab:extremos}
\resizebox{0.94\textwidth}{!}{\begin{tabular}{lrrrrr}
\toprule
 & Minimo & P1 & P99 & Maximo & Fuera de IQR \\
variable &  &  &  &  &  \\
\midrule
duration\_ms & 8586.000 & 61871.850 & 546010.980 & 5237295.000 & 4225 \\
tempo & 0.000 & 65.165 & 193.652 & 243.372 & 514 \\
loudness & -49.531 & -28.542 & -1.628 & 4.532 & 5026 \\
speechiness & 0.000 & 0.025 & 0.619 & 0.965 & 10644 \\
instrumentalness & 0.000 & 0.000 & 0.958 & 1.000 & 19613 \\
liveness & 0.000 & 0.041 & 0.950 & 1.000 & 6981 \\
\bottomrule
\end{tabular}
}
\end{table}

\begin{figure}[H]
\centering
\includegraphics[width=0.93\textwidth]{10_outlier_review.png}
\caption{Revisión de extremos. Las colas largas y acumulaciones en cero requieren interpretación por
variable, no eliminación automática.}
\label{fig:revision-extremos}
\end{figure}

\code{duration_ms} alcanza 5.237.295 ms (87,29 minutos); puede representar contenido largo válido.
\code{tempo} incluye cero y merece validación puntual, pero no se infiere que sea inválido sin
metadatos adicionales. \code{loudness} llega a $-49{,}531$ y 4,532 dB. En
\code{speechiness}, \code{instrumentalness} y \code{liveness}, una gran cantidad fuera de IQR es
compatible con distribuciones asimétricas acotadas 0--1. La implicancia es probar transformaciones o
escaladores robustos y analizar sensibilidad; eliminar puntos podría borrar géneros completos o
casos legítimos.
